
\subsection*{7.13. Relations avec la cohomologie @@M0@@-adique}
Rappelons (SGA 5 VII) que pour tout Module localement libre de type fini @@M1@@ sur un schéma @@M2@@, et tout nombre premier distinct des caractéristiques résiduelles de @@M3@@, on définit des classes de Chern @@M4@@-adiques (pour @@M5@@)
@@M6@@
satisfaisant à une formule d'additivité
@@M7@@
lorsque @@M8@@ est extension d'un module localement libre @@M9@@ par un module localement libre @@M10@@. De cette additivité on conclut que la classe de Chern totale @@M11@@ provient d'un homomorphisme
@@M12@@
(où l'objet but est un groupe par le produit, comme expliqué dans V 1, et où le @@M13@@ est le @@M14@@ «~\emph{naïf}~»). On en déduit des applications canoniques
@@M15@@

Ceci posé, soit @@M16@@ un schéma \emph{propre} et plat sur @@M17@@, à \emph{fibres régulières}, et supposons @@M18@@ \emph{hensélien}. On suppose @@M19@@ car. résiduelle de @@M20@@. On peut alors former le diagramme
@@M21@@
où les flèches horizontales sont les applications «~classe de Chern~» relatives à @@M22@@ et @@M23@@ respectivement, et où @@M24@@ est l'homomorphisme composé
@@M25@@
